\documentclass[12pt]{article}

\usepackage{sbc-template}
\usepackage{graphicx,url}
\usepackage[utf8]{inputenc}
\usepackage[T1]{fontenc}
\usepackage[brazilian]{babel}
\usepackage{booktabs}
\usepackage{array}
\usepackage{multirow}
\IfFileExists{placeins.sty}{\usepackage{placeins}}{\newcommand{\FloatBarrier}{}}
\usepackage{tikz}
\usepackage{pgfplots}
\usepackage{pgfplotstable}
\pgfplotsset{compat=1.18}
\usepgfplotslibrary{groupplots}
\usetikzlibrary{positioning,arrows.meta}

\graphicspath{{figuras/}}
\urlstyle{same}
\sloppy

\title{Pseudo-Box Target Sensitivity for Point-Supervised Cervical Cell Detection and Counting on CRIC Cervix}

\author{Anonymous submission}

\address{Author names and affiliations omitted for review}

\begin{document}

\maketitle

\begin{abstract}
Point annotations identify cervical cell centers but do not define object
extent. This work studies the experimental consequences of converting such
points into detection targets on the CRIC Cervix dataset, across ten pseudo-box
sizes (48--192 pixels) and three YOLO11 scales with image-disjoint partitions.
Validation F1 forms a broad plateau between 96 and 192 pixels, with distinct
precision, recall and counting trade-offs; the 144-pixel target was kept as an
intermediate operating point and YOLO11s gave the best balance. On the held-out
test set it reached AP50=0.8723, F1=0.8190 and a mean absolute counting error of
4.68 cells per image. A point-distance criterion, independent of the synthetic
box, yields an equivalent F1, confirming genuine nucleus localization. The study
contributes an auditable target-construction analysis for point-supervised
cervical cell detection; it claims neither anatomical bounding boxes nor
diagnostic use.
\end{abstract}

\begin{resumo}
Anota\c{c}\~oes pontuais identificam centros celulares, mas n\~ao definem a
extens\~ao do objeto. Este trabalho estuda as consequ\^encias experimentais de
converter esses pontos em alvos de detec\c{c}\~ao na base CRIC Cervix, sobre dez
tamanhos de pseudo-caixa (48--192 pixels) e tr\^es escalas YOLO11, com
parti\c{c}\~oes separadas por imagem. O F1 de valida\c{c}\~ao forma um plat\^o
entre 96 e 192 pixels, com compromissos distintos de precis\~ao, revoca\c{c}\~ao e
contagem; o alvo de 144 pixels foi mantido como ponto operacional intermedi\'ario
e o YOLO11s deu o melhor equil\'ibrio. No teste retido, alcan\c{c}ou AP50=0,8723,
F1=0,8190 e erro absoluto m\'edio de 4,68 c\'elulas por imagem. Um crit\'erio de
dist\^ancia ao ponto, independente da caixa sint\'etica, produz F1 equivalente,
confirmando localiza\c{c}\~ao real dos n\'ucleos. A contribui\c{c}\~ao \'e uma
an\'alise audit\'avel da constru\c{c}\~ao do alvo para detec\c{c}\~ao cervical
supervisionada por pontos, sem reivindicar caixas anat\^omicas ou uso
diagn\'ostico.
\end{resumo}

\section{Introdu\c{c}\~ao}

O rastreamento do c\^ancer do colo do \'utero depende de programas capazes de
identificar altera\c{c}\~oes precursoras antes da progress\~ao da doen\c{c}a. Em
2020, foram estimados aproximadamente 604 mil novos casos e 342 mil \'obitos no
mundo, com maior carga onde vacina\c{c}\~ao, rastreamento e tratamento s\~ao menos
acess\'iveis \cite{singh2023global}. Na citologia convencional, a automa\c{c}\~ao
tem sido estudada em diferentes escalas: classifica\c{c}\~ao de recortes,
segmenta\c{c}\~ao de estruturas, detec\c{c}\~ao de inst\^ancias e an\'alise de campos
ou l\^aminas \cite{jiang2023systematic,fang2024systematic}.

Essas tarefas pressup\~oem formas distintas de anota\c{c}\~ao. Um r\'otulo de
classe informa o que h\'a no recorte; um ponto informa onde uma inst\^ancia est\'a;
uma caixa acrescenta extens\~ao aproximada; uma m\'ascara descreve fronteira.
Logo, transformar um ponto em caixa n\~ao \'e apenas convers\~ao de formato: \'e uma
decis\~ao sobre qual regi\~ao o detector deve aprender. Em campos com muitas
c\'elulas, o tamanho adotado tamb\'em controla a sobreposi\c{c}\~ao entre alvos e
altera a rela\c{c}\~ao entre localiza\c{c}\~ao e contagem.

A CRIC Cervix disponibiliza 400 imagens convencionais de Papanicolau e 11.534
c\'elulas classificadas, cada uma associada a uma coordenada nuclear
\cite{rezende2021cric}. A base permite estudar detec\c{c}\~ao sem recortar
previamente as c\'elulas, mas n\~ao fornece bounding boxes anat\^omicas. Trabalhos
anteriores usaram a CRIC em detec\c{c}\~ao e classifica\c{c}\~ao
\cite{kalbhor2023cervicell}; ainda assim, a sensibilidade do resultado \`a
geometria criada a partir dos pontos permanece uma quest\~ao experimental
n\~ao investigada. A contagem autom\'atica de c\'elulas por imagem \'e uma tarefa
complementar ao diagn\'ostico --- na contagem celular por redes de regress\~ao,
o erro absoluto m\'edio depende tanto do m\'etodo quanto da densidade dos campos
\cite{xie2018fcrn} --- e, como se ver\'a, o tamanho da pseudo-caixa afeta
localiza\c{c}\~ao e contagem de formas distintas.

Este artigo formula o problema como \emph{engenharia do alvo supervisionado}:
n\~ao se prop\~oe uma nova arquitetura nem se usa a categoria Bethesda como sa\'ida.
Todas as coordenadas s\~ao tratadas como inst\^ancias de uma classe \'unica e o
objetivo \'e determinar em que medida um detector treinado com pseudo-caixas
recupera os pontos anotados e preserva a contagem por campo.

Tr\^es perguntas orientam o desenho experimental. \textbf{RQ1 (alvo):} quanto o
tamanho da pseudo-caixa modifica F1, AP50 e erro de contagem?
\textbf{RQ2 (capacidade):} aumentar a escala de YOLO11n para YOLO11s ou YOLO11m
melhora o compromisso entre localiza\c{c}\~ao e contagem?
\textbf{RQ3 (transfer\^encia):} a configura\c{c}\~ao selecionada na valida\c{c}\~ao
mant\'em desempenho no conjunto de teste retido, e como o erro varia com a
densidade celular?

A principal contribui\c{c}\~ao \'e tornar mensur\'avel uma decis\~ao normalmente
escondida na prepara\c{c}\~ao dos dados: a geometria do alvo criada a partir de
pontos. Todos os resultados s\~ao acompanhados de parti\c{c}\~oes por imagem,
limiar escolhido sem acesso ao teste e intervalos de confian\c{c}a calculados por
bootstrap sobre imagens completas.

\FloatBarrier
\section{Supervis\~ao Espacial em Imagens Celulares}

\subsection{Do contorno completo ao ponto}

M\'ascaras de inst\^ancia permitem avaliar fronteiras e separar objetos
adjacentes, mas exigem anota\c{c}\~ao detalhada. HoVer-Net, por exemplo, aprende
segmenta\c{c}\~ao e classifica\c{c}\~ao nuclear em histologia a partir de contornos
de inst\^ancia \cite{graham2019hovernet}; Cellpose tamb\'em aborda segmenta\c{c}\~ao
celular generalista com supervis\~ao de m\'ascaras \cite{stringer2021cellpose}.
Esses m\'etodos resolvem um problema espacial mais rico que o deste artigo, mas
dependem de um tipo de ground truth ausente na CRIC.

Pontos reduzem custo de anota\c{c}\~ao e identificam inst\^ancias sem descrever
suas fronteiras. Ribera et al. estudaram localiza\c{c}\~ao de objetos sem bounding
boxes \cite{ribera2019locating}, e Cheng et al. investigaram segmenta\c{c}\~ao de
inst\^ancias supervisionada por pontos \cite{cheng2022pointly}. Em imagens
biom\'edicas, Chen et al. propuseram segmenta\c{c}\~ao histopatol\'ogica fracamente
supervisionada com pontos esparsos \cite{chen2021sparsepoints}; Xia et al.
exploraram pontos e coloriza\c{c}\~ao para segmenta\c{c}\~ao nuclear
\cite{xia2023colorization}. Em comum, esses estudos tratam o ponto como supervis\~ao incompleta, n\~ao como
caixa impl\'icita. Uma abordagem alternativa \'e aprender a extens\~ao diretamente
por regress\~ao, como faz o P2BNet \cite{chen2023p2bnet}. O presente trabalho,
ao contr\'ario, torna \emph{expl\'icita} a geometria fixa e mede seu efeito ---
uma escolha que m\'etodos como o P2BNet ocultam na otimiza\c{c}\~ao.

Neste artigo, pseudo-caixas s\~ao adotadas por pragmatismo: permitem usar um
detector convencional mantendo sa\'ida por inst\^ancia. O IoU mede concord\^ancia
com um alvo sint\'etico, n\~ao com a fronteira citol\'ogica; por isso, o tamanho
do alvo \'e submetido a abla\c{c}\~ao e a contagem avaliada em paralelo.

\subsection{Detec\c{c}\~ao em citologia cervical}

Revis\~oes da \'area descrevem resultados elevados em classifica\c{c}\~ao celular,
mas tamb\'em heterogeneidade de bases, prepara\c{c}\~ao, classes e protocolos de
avalia\c{c}\~ao \cite{jiang2023systematic,fang2024systematic}. Essa heterogeneidade
impede interpretar m\'etricas de tarefas diferentes como ranking direto.
Classificar um recorte centrado, localizar uma c\'elula no campo e segmentar
n\'ucleo/citoplasma n\~ao s\~ao problemas intercambi\'aveis.

O CerviCell-detector avaliou detectores de objetos em imagens de Papanicolau,
incluindo experimentos na CRIC, com formula\c{c}\~oes bin\'aria e Bethesda
\cite{kalbhor2023cervicell}. O presente estudo remove a decis\~ao diagn\'ostica do
alvo e pergunta se todas as inst\^ancias pontuais podem ser localizadas e
contadas. Essa escolha produz uma contribui\c{c}\~ao complementar: um detector de
classe \'unica pode anteceder um classificador, mas seu desempenho deve primeiro
ser entendido como problema de cobertura espacial.

AP em diferentes limiares de IoU \'e adotada para medir rigor de localiza\c{c}\~ao,
seguindo a tradi\c{c}\~ao de benchmarks como COCO \cite{lin2014coco}. Entretanto,
como o alvo \'e sint\'etico, AP n\~ao \'e analisada isoladamente. F1 em um ponto
operacional e erros de contagem completam a avalia\c{c}\~ao, enquanto o conjunto
de teste permanece oculto durante a sele\c{c}\~ao.

\FloatBarrier
\section{Protocolo Experimental}

\subsection{Dados, unidade amostral e parti\c{c}\~oes}

Foram usadas as 400 imagens e as 11.534 coordenadas celulares da CRIC Cervix
\cite{rezende2021cric}. As categorias Bethesda foram preservadas nos metadados
para auditoria, mas n\~ao participaram do treinamento: a classe \'unica
\textit{cell} agrega as seis categorias da base, de c\'elulas normais (58,8\%) a
les\~oes de alto grau e carcinoma (HSIL 14,8\%, LSIL 11,8\%, ASC-H 8,0\%, ASC-US
5,3\% e SCC 1,4\%). Assim, a tarefa \'e de localiza\c{c}\~ao agn\'ostica \`a classe
sobre um conjunto clinicamente heterog\^eneo. Cada coordenada gerou uma inst\^ancia
dessa classe.

A imagem foi a unidade de parti\c{c}\~ao, na propor\c{c}\~ao 80/10/10. Para
evitar que a carga celular se concentrasse em uma parti\c{c}\~ao, as imagens
foram ordenadas pelo n\'umero de c\'elulas e atribu\'idas de forma gulosa \`a
parti\c{c}\~ao mais distante de sua meta proporcional de objetos (semente 42). O
procedimento destinou 320 imagens ao treino, 40 \`a valida\c{c}\~ao e 40 ao teste,
com 9.236, 1.149 e 1.149 c\'elulas, respectivamente; o balanceamento explica a
igualdade exata de c\'elulas entre valida\c{c}\~ao e teste. Nenhuma imagem foi
compartilhada entre parti\c{c}\~oes. A CRIC n\~ao disponibiliza identificador de
paciente; logo, a separa\c{c}\~ao por imagem \'e o maior controle de independ\^encia
suportado pelos metadados, n\~ao uma garantia de independ\^encia cl\'inica.

\subsection{Constru\c{c}\~ao do alvo}

Para um ponto nuclear $(x_i,y_i)$ e lado $s$, definiu-se a pseudo-caixa

\begin{equation}
B_i(s)=\left[x_i-\frac{s}{2},\ y_i-\frac{s}{2},\ x_i+\frac{s}{2},\ y_i+\frac{s}{2}\right],
\label{eq:pseudobox}
\end{equation}

recortada aos limites da imagem. Centro, largura e altura foram normalizados no
formato YOLO. A transforma\c{c}\~ao fixa o centro fornecido pelo especialista,
mas introduz uma hip\'otese sobre extens\~ao. Essa hip\'otese foi testada para
$s\in\{48,64,72,96,112,128,144,160,176,192\}$ pixels.

A Figura~\ref{fig:pseudo-boxes} sobrep\~oe as pseudo-caixas geradas a uma imagem
real da CRIC, permitindo inspecionar visualmente o alinhamento entre o alvo
sint\'etico e a citologia subjacente.

\begin{figure}[htbp]
\centering
\includegraphics[width=.56\textwidth]{pseudo_boxes.png}
\caption{Auditoria visual de pseudo-caixas sobre imagem real da CRIC.
Caixas s\~ao alvos computacionais, n\~ao contornos anat\^omicos.}
\label{fig:pseudo-boxes}
\end{figure}

\subsection{Fases do experimento}

A Figura~\ref{fig:desenho} separa constru\c{c}\~ao, sele\c{c}\~ao e teste. Na fase
RQ1, YOLO11n \cite{khanam2024yolov11} foi treinado para cada valor de $s$, por
at\'e 35 \'epocas, para isolar a sensibilidade ao alvo. Na RQ2, fixou-se $s=144$
e foram comparados YOLO11n, YOLO11s e YOLO11m. A RQ3 usou apenas o YOLO11s
selecionado e o limiar definido na valida\c{c}\~ao para abrir o teste uma \'unica
vez.

\begin{figure}[htbp]
\centering
\begin{tikzpicture}[
  node distance=7mm and 6mm,
  box/.style={draw=black!65,align=center,minimum height=7mm,text width=2.65cm,font=\scriptsize},
  arrow/.style={-{Latex[length=2mm]},thick,draw=black!65}
]
\node[box] (points) {CRIC\\pontos nucleares};
\node[box,right=of points] (targets) {10 tamanhos\\de pseudo-caixa};
\node[box,right=of targets] (rq1) {RQ1\\sensibilidade do alvo};
\node[box,below=of rq1] (rq2) {RQ2\\YOLO11n/s/m};
\node[box,left=of rq2] (select) {Valida\c{c}\~ao\\modelo e limiar};
\node[box,left=of select] (test) {RQ3\\teste retido};
\draw[arrow] (points)--(targets);
\draw[arrow] (targets)--(rq1);
\draw[arrow] (rq1)--(rq2);
\draw[arrow] (rq2)--(select);
\draw[arrow] (select)--(test);
\end{tikzpicture}
\caption{Fluxo experimental. O teste n\~ao participa da escolha de caixa,
capacidade ou limiar.}
\label{fig:desenho}
\end{figure}

\FloatBarrier
Os modelos foram inicializados com pesos pr\'e-treinados. Usou-se AdamW, taxa
inicial 0,002, agendamento cossenoidal, \textit{mosaic} moderado, varia\c{c}\~oes
geom\'etricas/crom\'aticas e precis\~ao mista. YOLO11n e YOLO11s receberam entrada
nominal 832; YOLO11m, 864. Os lotes foram 12, 6 e 4, respectivamente. As
compara\c{c}\~oes de capacidade tiveram 80 \'epocas m\'aximas e paci\^encia 20. O
treinamento foi executado em uma RTX 3060 Laptop com 6 GB de VRAM.

\subsection{Crit\'erios de avalia\c{c}\~ao}

As predi\c{c}\~oes foram exportadas pela rotina de valida\c{c}\~ao do Ultralytics,
sem aumento no teste, e todas as m\'etricas foram recomputadas a partir delas. A
associa\c{c}\~ao predi\c{c}\~ao--alvo priorizou maior IoU e imp\^os unicidade de
ambos os lados. A AP em IoU 0,30, 0,50 e 0,75 foi obtida por interpola\c{c}\~ao de
101 pontos, no estilo COCO \cite{lin2014coco}; precis\~ao, revoca\c{c}\~ao e F1
completam o conjunto. Como a pseudo-caixa \'e sint\'etica, adotou-se ainda um
crit\'erio independente da sua geometria: uma detec\c{c}\~ao \'e correta se o centro
da predi\c{c}\~ao estiver a at\'e $R$ pixels do ponto nuclear anotado, com
associa\c{c}\~ao gulosa \'unica por menor dist\^ancia. Esse crit\'erio avalia o
detector contra a anota\c{c}\~ao original, n\~ao contra a caixa derivada dela. O limiar de confian\c{c}a foi escolhido, em cada configura\c{c}\~ao, pelo maior
F1 na valida\c{c}\~ao em IoU 0,50. Na fase RQ1, o valor \'otimo do YOLO11n
variou entre 0,20 e 0,50 conforme o tamanho da pseudo-caixa, refletindo a
calibra\c{c}\~ao individual de cada treino. Na fase RQ2/RQ3, os tr\^es modelos
comparados (YOLO11n, YOLO11s e YOLO11m) atingiram o pico em 0,35; o valor do
YOLO11s foi usado na abertura do teste.

Para cada imagem $j$, com contagem anotada $g_j$ e predita $p_j$, foram
calculados

\begin{equation}
\mathrm{MAE}=\frac{1}{N}\sum_j|p_j-g_j|,\qquad
\mathrm{bias}=\frac{1}{N}\sum_j(p_j-g_j).
\end{equation}

RMSE e correla\c{c}\~ao complementaram a an\'alise. Intervalos de 95\% foram
obtidos com 1.000 r\'eplicas de bootstrap sobre imagens completas. Essa unidade
de reamostragem preserva a depend\^encia entre c\'elulas do mesmo campo.

\subsection{Rastreabilidade e controle de decis\~oes}

Cada etapa produziu artefatos persistentes: lista de imagens por parti\c{c}\~ao,
metadados das pseudo-caixas, configura\c{c}\~ao de treino, hist\'orico por \'epoca,
pesos selecionados, predi\c{c}\~oes por imagem e tabelas de m\'etricas. Figuras e
valores reportados derivam diretamente desses arquivos, sem transcri\c{c}\~ao
manual. A origem dos dados, o particionamento, o hardware, a sele\c{c}\~ao
do limiar e as limita\c{c}\~oes de generaliza\c{c}\~ao foram explicitados em
conson\^ancia com recomenda\c{c}\~oes de relato para IA em imagens m\'edicas
\cite{mongan2020claim,tejani2024claim,collins2024tripod}.

As decis\~oes foram encadeadas de modo unilateral: RQ1 definiu a regi\~ao de
pseudo-caixas; RQ2 comparou capacidades com o mesmo alvo; a valida\c{c}\~ao
selecionou confian\c{c}a; e somente ent\~ao o teste foi avaliado. Nenhuma m\'etrica
de teste foi usada para alterar pesos, lado da caixa ou limiar. Os resultados do
sweep, por sua vez, s\~ao tratados como an\'alise explorat\'oria de sensibilidade,
pois dez configura\c{c}\~oes e uma execu\c{c}\~ao por configura\c{c}\~ao n\~ao
permitem atribuir toda diferen\c{c}a observada exclusivamente ao tamanho do alvo.

\FloatBarrier
\section{Resultados}

\subsection{RQ1: o alvo altera o resultado?}

A Figura~\ref{fig:sweep-pgf} mostra que a resposta \'e n\~ao monot\^onica. F1 e
AP50 aumentaram das menores caixas at\'e uma regi\~ao de estabilidade, mas o
tamanho 160 apresentou queda acentuada e supercontagem. O maior F1 foi 0,8253
em 192 pixels. O menor MAE (3,65) e o maior AP50 (0,8776) ocorreram em 176
pixels. Em 144 pixels, F1 foi 0,8228, AP50 0,8748 e MAE 4,40.

\begin{figure}[htbp]
\centering
\begin{tikzpicture}
\begin{groupplot}[
 group style={group size=2 by 1,horizontal sep=1.25cm},
 width=.46\textwidth,height=5.1cm,
 xlabel={Lado da pseudo-caixa (px)},
 xmin=40,xmax=200,xtick={48,72,96,120,144,168,192},
 tick label style={font=\scriptsize},label style={font=\small},
 grid=major,grid style={black!10},axis line style={black!65}
]
\nextgroupplot[ylabel={Desempenho},ymin=.62,ymax=.96,legend style={font=\scriptsize,draw=none,at={(.5,-.34)},anchor=north,legend columns=3}]
\addplot[black,thick,mark=*] table[x=box_size_px,y=f1_val_iou50,col sep=comma]{dados/bbox_sweep.csv};
\addlegendentry{F1@0,50}
\addplot[black!50,thick,dashed,mark=square*] table[x=box_size_px,y=ap50_val,col sep=comma]{dados/bbox_sweep.csv};
\addlegendentry{AP50}
\addplot[black!30,thick,densely dotted,mark=triangle*] table[x=box_size_px,y=recall_val_iou50,col sep=comma]{dados/bbox_sweep.csv};
\addlegendentry{Revoca\c{c}\~ao}
\addplot[only marks,mark=*,mark size=3pt] coordinates {(144,0.8228176)};
\nextgroupplot[ylabel={MAE de contagem},ymin=0,ymax=29]
\addplot[black!55,thick,mark=triangle*] table[x=box_size_px,y=mae_contagem_val,col sep=comma]{dados/bbox_sweep.csv};
\addplot[only marks,mark=triangle*,mark size=3pt] coordinates {(144,4.4)};
\end{groupplot}
\end{tikzpicture}
\caption{Sensibilidade ao tamanho do alvo na valida\c{c}\~ao. O marcador maior
indica a configura\c{c}\~ao de 144 pixels usada na compara\c{c}\~ao de capacidade.}
\label{fig:sweep-pgf}
\end{figure}

Os valores entre 96 e 192 pixels, exceto 160, formam um plat\^o estreito de F1
(0,8157--0,8253). Assim, os dados n\~ao sustentam um \'unico tamanho universal.
O ponto de 144 pixels foi mantido por combinar AP50 pr\'oxima do m\'aximo
(0,8748 vs.\ 0,8776 em 176), revoca\c{c}\~ao superior a 176 e 192
(0,8285 vs.\ 0,7885 e 0,8059) e posi\c{c}\~ao intermedi\'aria no plat\^o. A maior
revoca\c{c}\~ao importa quando as c\'elulas perdidas pelo detector representam
potenciais falsos negativos diagn\'osticos; pagar um custo de MAE levemente
maior (4,40 vs.\ 3,65 em 176) para manter cobertura foi o crit\'erio adotado.
A instabilidade em 160 pixels impede interpretar a tend\^encia como simples
benef\'icio de caixas maiores. Naquele ponto, o detector saturou o teto de
detec\c{c}\~oes (300 por imagem em todas as 40 imagens de valida\c{c}\~ao), com
revoca\c{c}\~ao 0,955 e precis\~ao 0,496. Como caixas \emph{maiores} (176 e 192
pixels) n\~ao reproduzem esse colapso, o efeito reflete a vari\^ancia de um treino
isolado, n\~ao uma consequ\^encia geom\'etrica do tamanho do alvo. De fato, a
mAP50 nativa desse modelo permaneceu saud\'avel (0,883), indicando que o ranking
das detec\c{c}\~oes foi preservado e que o colapso se restringiu \`a calibra\c{c}\~ao
do ponto operacional.

\FloatBarrier
\subsection{RQ2: maior capacidade melhora a tarefa?}

YOLO11s apresentou o maior F1 (0,8363), AP50 (0,8848) e o menor MAE (4,68) na
valida\c{c}\~ao. YOLO11n ficou pr\'oximo, com F1 0,8231 e MAE 4,83. YOLO11m
atingiu a maior revoca\c{c}\~ao (0,8921), mas sua precis\~ao de 0,7158 elevou o
MAE para 8,13 e reduziu F1 para 0,7943. A Figura~\ref{fig:modelos-pgf} mostra
que escala de par\^ametros e desempenho operacional n\~ao cresceram juntos.

\begin{figure}[htbp]
\centering
\begin{tikzpicture}
\begin{axis}[
 width=.76\textwidth,height=5.2cm,
 ybar,bar width=7pt,
 symbolic x coords={YOLO11n,YOLO11s,YOLO11m},xtick=data,
 ymin=.68,ymax=.92,ylabel={M\'etrica na valida\c{c}\~ao},
 tick label style={font=\small},label style={font=\small},
 legend style={font=\scriptsize,draw=none,at={(.5,-.22)},anchor=north,legend columns=3},
 ymajorgrids=true,grid style={black!10},axis line style={black!65}
]
\addplot[fill=black!75,draw=black!75] table[x=modelo,y=f1_val_iou50,col sep=comma]{dados/modelos.csv};
\addlegendentry{F1}
\addplot[fill=black!45,draw=black!55] table[x=modelo,y=revocacao_val_iou50,col sep=comma]{dados/modelos.csv};
\addlegendentry{Revoca\c{c}\~ao}
\addplot[fill=black!15,draw=black!45] table[x=modelo,y=ap50_val_custom,col sep=comma]{dados/modelos.csv};
\addlegendentry{AP50}
\end{axis}
\end{tikzpicture}
\caption{Compara\c{c}\~ao das escalas YOLO11 com pseudo-caixas de 144 pixels.}
\label{fig:modelos-pgf}
\end{figure}

O ganho de YOLO11s sobre YOLO11n foi de 0,0132 em F1 e 0,0075 em AP50. Essa
diferen\c{c}a \'e pequena e n\~ao possui intervalo de confian\c{c}a entre treinos;
portanto, suporta a escolha operacional do S neste experimento, mas n\~ao uma
afirma\c{c}\~ao geral de superioridade arquitetural. O M foi descartado porque
o aumento de revoca\c{c}\~ao n\~ao compensou falsos positivos e erro de contagem.

\FloatBarrier
\subsection{Sele\c{c}\~ao do ponto operacional}

A confian\c{c}a controla uma troca relevante para contagem. Na valida\c{c}\~ao,
limiares baixos preservaram quase todas as c\'elulas, mas geraram muitas
detec\c{c}\~oes excedentes: em 0,10, a revoca\c{c}\~ao foi 0,9661 e a precis\~ao
0,5187, com MAE 24,78. Em 0,50, a precis\~ao subiu para 0,8853, mas a revoca\c{c}\~ao
caiu para 0,6919 e o vi\'es tornou-se negativo ($-6,28$ c\'elulas por imagem).

\begin{figure}[htbp]
\centering
\begin{tikzpicture}
\begin{axis}[
 width=.72\textwidth,height=5.4cm,
 xlabel={Limiar de confian\c{c}a},ylabel={F1 na valida\c{c}\~ao},
 xmin=.02,xmax=.52,ymin=.40,ymax=.88,
 xtick={.03,.10,.20,.25,.35,.50},
 tick label style={font=\scriptsize},label style={font=\small},
 legend style={font=\scriptsize,draw=none,at={(.5,-.40)},anchor=north,legend columns=3},
 grid=major,grid style={black!10},axis line style={black!65}
]
\addplot[black!75,thick,mark=*] table[x=conf,y=f1_iou30,col sep=comma]{dados/f1_limiares.csv};
\addlegendentry{IoU 0,30}
\addplot[black!45,thick,dashed,mark=square*] table[x=conf,y=f1_iou50,col sep=comma]{dados/f1_limiares.csv};
\addlegendentry{IoU 0,50}
\addplot[black!20,thick,dotted,mark=triangle*] table[x=conf,y=f1_iou75,col sep=comma]{dados/f1_limiares.csv};
\addlegendentry{IoU 0,75}
\addplot[only marks,mark=*,mark size=3.2pt] coordinates {(.35,.8363034)};
\end{axis}
\end{tikzpicture}
\caption{Curvas de F1 usadas para selecionar a confian\c{c}a. O marcador maior
identifica 0,35, escolhido exclusivamente na valida\c{c}\~ao em IoU 0,50.}
\label{fig:limiar-pgf}
\end{figure}

O limiar 0,35 maximizou F1@0,50 (0,8363), com precis\~ao 0,7983, revoca\c{c}\~ao
0,8782, MAE 4,68 e vi\'es 2,88. A proximidade das curvas em IoU 0,30 e 0,50
indica que, nesse intervalo, a maioria das associa\c{c}\~oes aceitas em 0,30
tamb\'em satisfaz 0,50. A separa\c{c}\~ao em IoU 0,75 evidencia a penaliza\c{c}\~ao
adicional por desalinhamento fino com a pseudo-caixa. O teste recebeu 0,35 sem
novo ajuste; portanto, seu resultado mede transfer\^encia do ponto operacional,
n\~ao o melhor limiar retrospectivo para aquela parti\c{c}\~ao.

\FloatBarrier
\subsection{RQ3: o resultado se mant\'em no teste?}

Com YOLO11s, lado 144 e confian\c{c}a 0,35, o teste produziu 975 verdadeiros
positivos, 257 falsos positivos e 174 falsos negativos. A
Tabela~\ref{tab:teste-unificado} re\'une todas as medidas, cada uma com intervalo
de 95\% obtido por bootstrap de imagens.

\begin{table}[htbp]
\centering
\caption{Resultado no teste retido. IC95\% por bootstrap de imagens.}
\label{tab:teste-unificado}
\small
\setlength{\tabcolsep}{5pt}
\begin{tabular}{lrl@{\hspace{8mm}}lrl}
\toprule
M\'etrica & Valor & IC95\% & M\'etrica & Valor & IC95\% \\
\midrule
Precis\~ao & 0,7914 & [0,746; 0,833] & AP30 & 0,8843 & [0,858; 0,913] \\
Revoca\c{c}\~ao & 0,8486 & [0,822; 0,877] & AP50 & 0,8723 & [0,843; 0,904] \\
F1 & 0,8190 & [0,796; 0,842] & AP75 & 0,7515 & [0,695; 0,811] \\
MAE & 4,675 & [3,150; 6,425] & RMSE & 7,216 & [4,560; 9,850] \\
Vi\'es & 2,075 & [0,050; 4,126] & Correla\c{c}\~ao & 0,9434 & [0,874; 0,979] \\
\bottomrule
\end{tabular}
\end{table}

O F1 de teste ficou 0,0173 abaixo da valida\c{c}\~ao, enquanto AP50 caiu 0,0125.
Essas diferen\c{c}as s\~ao compat\'iveis com transfer\^encia moderada entre duas
amostras de 40 imagens, sem evid\^encia de colapso. O vi\'es positivo indica 2,08
detec\c{c}\~oes excedentes por imagem em m\'edia.

\begin{figure}[htbp]
\centering
\begin{tikzpicture}
\begin{axis}[
 width=.66\textwidth,height=6.0cm,
 xlabel={C\'elulas anotadas por imagem},ylabel={C\'elulas detectadas por imagem},
 xmin=0,xmax=110,ymin=0,ymax=110,
 tick label style={font=\scriptsize},label style={font=\small},
 grid=major,grid style={black!10},axis line style={black!65}
]
\addplot[only marks,mark=*,mark size=1.7pt,black!60] table[x=gt_count,y=pred_count,col sep=comma]{dados/contagem_teste.csv};
\addplot[black,dashed,domain=0:110,samples=2]{x};
\end{axis}
\end{tikzpicture}
\caption{Concord\^ancia de contagem no teste; a linha tracejada representa
$p_j=g_j$.}
\label{fig:contagem-pgf}
\end{figure}

Na an\'alise explorat\'oria por tercis, o MAE foi 4,15 em imagens com 5--16
c\'elulas, 3,69 entre 17--30 e 6,07 entre 32--94. A faixa mais densa concentrou
maior erro absoluto, embora o F1 global permane\c{c}a 0,8190. A mediana global
do erro foi 3 c\'elulas. Esses estratos foram definidos depois da avalia\c{c}\~ao
principal e n\~ao devem ser interpretados como hip\'oteses confirmat\'orias.

\FloatBarrier
\subsection{Avalia\c{c}\~ao independente da pseudo-caixa}

Para verificar se o desempenho reflete localiza\c{c}\~ao real e n\~ao apenas
ader\^encia \`a caixa sint\'etica, o teste foi reavaliado pelo crit\'erio de
dist\^ancia ao ponto. Como n\'ucleos vizinhos distam em mediana 72 pixels, um raio
$R$ bem abaixo desse valor associa cada predi\c{c}\~ao ao n\'ucleo correto, n\~ao a
um adjacente. A Figura~\ref{fig:dist-pgf} mostra o F1 crescendo e estabilizando:
0,812 em $R=32$ pixels (metade do espa\c{c}amento t\'ipico) e 0,821 em 48.

\begin{figure}[htbp]
\centering
\begin{tikzpicture}
\begin{axis}[
 width=.70\textwidth,height=4.4cm,
 xlabel={Raio de toler\^ancia $R$ (px)},ylabel={M\'etrica no teste},
 xmin=6,xmax=50,ymin=.45,ymax=.92,
 xtick={8,16,24,32,40,48},
 tick label style={font=\scriptsize},label style={font=\small},
 legend style={font=\scriptsize,draw=none,at={(.5,-.30)},anchor=north,legend columns=3},
 grid=major,grid style={black!10},axis line style={black!65}
]
\addplot[black,thick,mark=*] table[x=R,y=f1,col sep=comma]{dados/dist_eval.csv};
\addlegendentry{F1}
\addplot[black!50,thick,dashed,mark=square*] table[x=R,y=precision,col sep=comma]{dados/dist_eval.csv};
\addlegendentry{Precis\~ao}
\addplot[black!25,thick,dotted,mark=triangle*] table[x=R,y=recall,col sep=comma]{dados/dist_eval.csv};
\addlegendentry{Revoca\c{c}\~ao}
\addplot[gray,densely dashdotted,domain=6:50,samples=2]{0.819};
\addlegendentry{F1 IoU 0,50}
\end{axis}
\end{tikzpicture}
\caption{Avalia\c{c}\~ao por dist\^ancia ao ponto anotado, independente da geometria
da pseudo-caixa. A linha horizontal \'e o F1 obtido por IoU 0,50.}
\label{fig:dist-pgf}
\end{figure}

Esse intervalo (0,812--0,821) coincide com o F1 por IoU 0,50 (0,819): duas medidas
de natureza distinta concordam, indicando que a qualidade reportada \'e
localiza\c{c}\~ao efetiva dos n\'ucleos, n\~ao um artefato da conven\c{c}\~ao de alvo.

\FloatBarrier
\section{Discuss\~ao}

\subsection{O que a pseudo-caixa representa}

O experimento n\~ao recupera fronteiras celulares; ele aprende uma vizinhan\c{c}a
quadrada ao redor do n\'ucleo. Essa distin\c{c}\~ao muda o significado das
m\'etricas. AP75 alta indica concord\^ancia com a regra da
Equa\c{c}\~ao~\ref{eq:pseudobox}, n\~ao segmenta\c{c}\~ao precisa de n\'ucleo ou
citoplasma. M\'etodos com m\'ascaras, como HoVer-Net e Cellpose, respondem a uma
pergunta espacial mais detalhada \cite{graham2019hovernet,stringer2021cellpose}.
A objec\c{c}\~ao de que o IoU contra um alvo sint\'etico seria igualmente
artificial \'e respondida pela avalia\c{c}\~ao por dist\^ancia: como o F1 por ponto
converge ao F1 por IoU, a m\'etrica n\~ao depende criticamente da geometria
adotada, e o n\'ucleo anotado --- \'unico ground truth da CRIC --- \'e de fato
recuperado.

Ainda assim, pontos cont\^em informa\c{c}\~ao de inst\^ancia que r\'otulos globais
n\~ao possuem. A literatura de supervis\~ao pontual mostra que essa informa\c{c}\~ao
pode orientar localiza\c{c}\~ao e segmenta\c{c}\~ao com menor custo de anota\c{c}\~ao
\cite{ribera2019locating,chen2021sparsepoints,xia2023colorization}. O resultado
da RQ1 acrescenta uma observa\c{c}\~ao pr\'atica: quando pontos s\~ao convertidos em
caixas, a geometria escolhida deve ser reportada e testada. Caso contr\'ario, uma
parte relevante do desempenho permanece escondida na prepara\c{c}\~ao do dado.

O plat\^o de 96--192 pixels \'e mais informativo que a escolha isolada de 144.
Ele mostra que diversas extens\~oes produzem F1 parecido, mas n\~ao o mesmo
compromisso de precis\~ao, revoca\c{c}\~ao e contagem. A queda em 160 pixels alerta
contra uma leitura determinista da curva: cada ponto corresponde a um treino e
inclui variabilidade de otimiza\c{c}\~ao. Repeti\c{c}\~oes por semente seriam
necess\'arias para separar efeito geom\'etrico de varia\c{c}\~ao estoc\'astica.

\subsection{Capacidade e custo computacional}

YOLO11s foi o melhor compromisso observado, mas sua vantagem sobre YOLO11n foi
pequena. Sem repeti\c{c}\~oes independentes, n\~ao h\'a base estat\'istica para
afirmar que o S \'e intrinsecamente superior. A conclus\~ao suportada \'e mais
restrita: sob os pesos, parti\c{c}\~ao e or\c{c}amento executados, S gerou o melhor
ponto operacional.

O M ilustra por que mAP de treinamento n\~ao deve ser o \'unico crit\'erio. Sua
revoca\c{c}\~ao foi a maior, mas o aumento de falsos positivos deteriorou contagem
e F1. Esse comportamento importa em campos densos: detectar mais candidatos
pode ser desej\'avel em uma etapa de busca, mas transfere custo para supress\~ao
de duplicidades ou revis\~ao posterior. O artigo, portanto, n\~ao trata escala do
modelo como proxy autom\'atica de qualidade.

\subsection{Escopo da compara\c{c}\~ao com a literatura}

Os valores de AP deste artigo n\~ao constituem um ranking direto contra HoVer-Net
ou Cellpose: esses m\'etodos estimam m\'ascaras, foram avaliados em outros
conjuntos e respondem a uma defini\c{c}\~ao de alvo distinta
\cite{graham2019hovernet,stringer2021cellpose}. Tampouco se deve comparar apenas
o n\'umero final com trabalhos de classifica\c{c}\~ao Bethesda. O
CerviCell-detector, por exemplo, investiga detec\c{c}\~ao de c\'elulas cancerosas,
enquanto aqui toda coordenada celular pertence a uma classe bin\'aria
\cite{kalbhor2023cervicell}. M\'etrica, unidade de an\'alise, parti\c{c}\~ao e
sem\^antica do r\'otulo precisam coincidir para sustentar superioridade num\'erica;
as revis\~oes da \'area mostram justamente grande heterogeneidade nesses elementos
\cite{jiang2023systematic,fang2024systematic}.

A contribui\c{c}\~ao compar\'avel \'a literatura de supervis\~ao fraca \'e outra:
quantificar o efeito de transformar um ponto em uma regi\~ao de treinamento e
mostrar como essa escolha repercute simultaneamente em localiza\c{c}\~ao e
contagem. Ribera et al. aprendem localiza\c{c}\~ao sem caixas, enquanto Chen et
al. e Xia et al. exploram pontos para segmenta\c{c}\~ao
\cite{ribera2019locating,chen2021sparsepoints,xia2023colorization}. Este estudo
n\~ao substitui essas formula\c{c}\~oes; fornece um protocolo simples e audit\'avel
para o cen\'ario em que um detector de caixas \'e treinado a partir de coordenadas.
Assim, a novidade reivindicada \'e metodol\'ogica e experimental, n\~ao uma alegada
superioridade universal do YOLO11s.

\subsection{Contagem e fontes de erro}

A correla\c{c}\~ao 0,9434 confirma que o modelo acompanha a varia\c{c}\~ao de
carga celular entre imagens; o MAE de 4,68 e o vi\'es de 2,08 revelam, por\'em,
erro residual e tend\^encia de supercontagem. Para refer\^encia de ordem de grandeza, Xie et al.\ aplicaram redes de regress\~ao
totalmente convolucionais \`a contagem celular em microscopia e obtiveram MAE
competitivo em diferentes densidades \cite{xie2018fcrn}. O valor obtido aqui
(MAE~4,68) situa-se na mesma casa de magnitude, embora a compara\c{c}\~ao seja
apenas indicativa: os conjuntos de dados, a supervis\~ao e a formula\c{c}\~ao
diferem. Em termos absolutos, a faixa mais
densa concentrou o maior erro (MAE~6,31): pequenas taxas de duplicidade ou
omiss\~ao acumulam-se quando o campo cont\'em dezenas de c\'elulas.

H\'a pelo menos tr\^es explica\c{c}\~oes, n\~ao separadas por este experimento:
duplicidade em grupos sobrepostos, quando o IoU entre pseudo-caixas adjacentes
fica abaixo do limiar de supress\~ao e dois hits sobrevivem para o mesmo n\'ucleo;
estruturas n\~ao celulares confundidas com o alvo; e c\'elulas vis\'iveis sem ponto
correspondente na anota\c{c}\~ao. Distinguir essas causas exige revis\~ao cega de
falsos positivos por especialista. O colapso isolado em 160 pixels \'e de natureza
distinta: ali o detector saturou o teto de detec\c{c}\~oes por imagem, um efeito
de instabilidade de treino que n\~ao se repete nas demais escalas e n\~ao deve ser
confundido com a supercontagem residual do modelo final.

\subsection{Limita\c{c}\~oes e perspectivas}

Quatro limita\c{c}\~oes delimitam a interpreta\c{c}\~ao. \emph{Constru\c{c}\~ao:}
m\'etricas de IoU descrevem consist\^encia com um alvo sint\'etico, n\~ao com a
fronteira citol\'ogica. \emph{Interna:} cada configura\c{c}\~ao foi treinada uma
vez; repeti\c{c}\~oes por semente seriam necess\'arias para separar efeito de $s$
de ru\'ido de otimiza\c{c}\~ao. \emph{Externa:} avalia\c{c}\~ao em base \'unica,
separa\c{c}\~ao por imagem (n\~ao por paciente), sem teste em outro laborat\'orio
ou equipamento. \emph{Conclus\~ao:} os intervalos de bootstrap quantificam
varia\c{c}\~ao entre imagens, mas n\~ao incerteza de treinamento; diferen\c{c}as
no plat\^o de F1 n\~ao devem ser generalizadas sem repeti\c{c}\~oes.

A extens\~ao imediata \'e comparar pseudo-caixas a formulac\~oes que aprendem
extens\~ao a partir dos pontos \cite{chen2023p2bnet}, repetir o sweep com
m\'ultiplas sementes e anotar uma amostra com contornos manuais para calibrar a
avalia\c{c}\~ao atual.

\FloatBarrier
\section{Conclus\~ao}

Este estudo tratou a convers\~ao de pontos nucleares em caixas como vari\'avel
cient\'ifica. Para a RQ1, observou-se um plat\^o de F1 entre 96 e 192 pixels, com
compromissos distintos de AP50, revoca\c{c}\~ao e contagem. Para a RQ2, YOLO11s
foi o melhor ponto operacional, embora a vantagem sobre YOLO11n tenha sido
pequena e YOLO11m tenha apresentado maior revoca\c{c}\~ao. Para a RQ3, o modelo
selecionado obteve AP50 0,8723, F1 0,8190 e MAE 4,68 no teste retido; o erro
absoluto aumentou em imagens mais densas.

A conclus\~ao principal n\~ao \'e que 144 pixels ou YOLO11s sejam universais. O
resultado relevante \'e que a defini\c{c}\~ao espacial criada a partir de pontos
altera de forma mensur\'avel detec\c{c}\~ao e contagem e, portanto, precisa fazer
parte do protocolo reportado. O pipeline produzido oferece uma refer\^encia
reprodut\'ivel para esse tipo de auditoria na CRIC. Sua extens\~ao natural \'e
substituir a geometria fixa por supervis\~ao diretamente pontual e comparar com
uma amostra de contornos anotados.

\section*{Artefatos de Reprodu\c{c}\~ao}

Os CSVs usados pelos gr\'aficos, configura\c{c}\~oes de treinamento, checkpoints
selecionados e rotinas de avalia\c{c}\~ao est\~ao dispon\'iveis no reposit\'orio
p\'ublico do projeto. O reposit\'orio inclui as parti\c{c}\~oes por imagem (semente
42), os metadados das pseudo-caixas para cada valor de $s$ e os arquivos de
predi\c{c}\~ao do conjunto de teste, permitindo reproduzir as m\'etricas reportadas
sem re-treinamento.


\bibliographystyle{sbc}
\bibliography{sbc-template}

\end{document}
